\documentclass[11pt]{article}
\usepackage{amsmath, amssymb, amsthm, enumerate}
\usepackage[margin=1in, top=1.2in, bottom=1.2in]{geometry}
\usepackage{parskip}
\usepackage[colorlinks=true, linkcolor=blue, citecolor=blue, urlcolor=blue]{hyperref}
\usepackage{cleveref}
\theoremstyle{definition}
\newtheorem{definition}{Definition}[section]
\newtheorem{theorem}{Theorem}[section]
\newtheorem{lemma}{Lemma}[section]
\newtheorem{proposition}{Proposition}[section]
\newtheorem{corollary}{Corollary}[section]
\theoremstyle{remark}
\newtheorem{remark}{Remark}[section]
\newtheorem{example}{Example}[section]
\setlength{\parindent}{0pt}
\setlength{\parskip}{6pt}
\begin{document}

\begin{center}
{\Large\textbf{Techiniques,  Examples And Counterexamples}}\\
\vspace{0.3cm}
{\normalsize\today}
\end{center}
\vspace{0.5cm}


% --- Page 1 ---
\begin{itemize}
    \item \textbf{Modules $(M, R, +)$}
    \begin{proof}
        Verify the axioms.
    \end{proof}

    \item \textbf{Submodules}
    \begin{proof}
        Verify the axioms (Check subgroup and for any $r \in R$, $m \in M$, $rm \in M$).
    \end{proof}

    \item \textbf{Homomorphism}
    \begin{proof}
        Verify the axioms.
    \end{proof}
\end{itemize}

% --- Page 2 ---
% blank page

% --- Page 3 ---
\begin{example}\label{ex:modules}
    \begin{enumerate}
        \item Vector space $\mathbb{C}$ (Ring over module is a field).
        \item $(V, T)$ where $V$ is a vector space and $T$ is a linear transform. Ring is constructed to be $F[x]$ where
              \[
              f(x) \cdot v := f(T)(v).
              \]
              Note: $(V, T) \longmapsto F[x]$ module.
              For $\longmapsto$ we define $T(x) := x \cdot v$.
    \end{enumerate}
\end{example}

% --- Page 4 ---
\begin{itemize}
    \item \textbf{Vector subspaces}

    \item In $\mathbb{Z}$-modules say $M$
    \begin{definition}\label{def:submodule_z}
        $N$ is a submodule iff $(N, +) \leq (CM, +)$.
    \end{definition}

    \item In $R$ is a $R$-module
    \begin{definition}\label{def:submodule_ring}
        $N$ is a submodule iff $N$ is an ideal.
    \end{definition}

    \item If $M$ is a $F[x]$-module $\cong (V, T)$
    \begin{definition}\label{def:submodule_fx}
        $W$ is a submodule iff $T(CW) \subseteq W$, i.e. $W$ is invariant under $T$.
    \end{definition}
\end{itemize}

% --- Page 5 ---
\begin{tabular}{|l|l|l|}
    \hline
    \multicolumn{3}{|c|}{Finitely generated and cyclic modules} \\
    \hline
    \multirow{2}{*}{Module} & \multicolumn{2}{|c|}{f.g.} \\
    \cline{2-3}
    & $f.g.$ & cyclic \\
    \hline
    $(V, \mathcal{F})$ & $\dim V < +\infty$ & $\dim V = \{0, 1\}$ \\
    \hline
    \multirow{2}{*}{$G$ is abelian group treated like $\mathbb{Z}$-module} & \multicolumn{2}{|c|}{f.g. as group} \\
    \cline{2-3}
    & & cyclic as group \\
    \hline
    \multirow{2}{*}{$R$ comm ring and $M \triangleleft R$, $M$ is $R$-module} & \multicolumn{2}{|c|}{f.g. as an ideal} \\
    \cline{2-3}
    & & principal ideal \\
    \hline
\end{tabular}

\begin{definition}\label{def:finitely_generated_module}
    A module $(V, \mathcal{F})$ is called \emph{finitely generated} if $\dim V < +\infty$.
\end{definition}

\begin{definition}\label{def:cyclic_module}
    A module $(V, \mathcal{F})$ is called \emph{cyclic} if $\dim V = \{0, 1\}$.
\end{definition}

\begin{definition}\label{def:finitely_generated_abelian_group}
    An abelian group $G$, treated as a $\mathbb{Z}$-module, is called \emph{finitely generated as a group} if it is finitely generated as a $\mathbb{Z}$-module.
\end{definition}

\begin{definition}\label{def:cyclic_abelian_group}
    An abelian group $G$, treated as a $\mathbb{Z}$-module, is called \emph{cyclic as a group} if it is cyclic as a $\mathbb{Z}$-module.
\end{definition}

\begin{definition}\label{def:finitely_generated_ideal}
    Let $R$ be a commutative ring and $M \triangleleft R$ be an ideal of $R$. The ideal $M$ is called \emph{finitely generated as an ideal} if it is finitely generated as an $R$-module.
\end{definition}

\begin{definition}\label{def:principal_ideal}
    Let $R$ be a commutative ring and $M \triangleleft R$ be an ideal of $R$. The ideal $M$ is called a \emph{principal ideal} if it is cyclic as an $R$-module.
\end{definition}

% --- Page 6 ---
\begin{itemize}
    \item \begin{definition}\label{def:module_homomorphism}
        A \underline{module homomorphism} is a structure-preserving map between modules.
    \end{definition}

    \item Abelian groups $\leftrightarrow$ $\mathbb{Z}$-modules, then group homomorphism $\leftrightarrow$ module homomorphism.

    \item \begin{definition}\label{def:linear_transformation}
        For vector fields, a \underline{linear transformation} $\leftrightarrow$ module homomorphism.
    \end{definition}

    \item \begin{definition}\label{def:ring_homomorphism}
        For a ring $R$, a map $\Theta: M \to N$ where $M = N = R$ and $\Theta(rs) = r\Theta(s)$ is \textbf{not} a ring homomorphism but is a module homomorphism.
    \end{definition}

    \item \begin{remark}\label{rem:isomorphism_theorems}
        Isomorphism theorems hold for modules.
    \end{remark}
\end{itemize}

% --- Page 7 ---
% blank page

\end{document}